\clearpage
\section{Описание данных}
\label{sec:data}

\subsection{Источники и даты наблюдения}

Для исследования использованы три выгрузки каталога компаний Y Combinator с платформы Kaggle~\cite{kaggle-yc-dataset1,kaggle-yc-dataset2}. Файлы за февраль и июль 2023 года имеют одинаковые поля. Поздний файл относится к августовской версии 2026 года; в программе для неё используется обозначение 01.08.2026. Состав файлов и их роль приведены в Таблице~\ref{tab:data-snapshots}. Обозначение версии не означает, что каждое поле карточки было обновлено в день выпуска версии.

\begin{table}[!htbp]
    \centering
    \small
    \caption{Исходные файлы и их роль в исследовании. Число строк указано до удаления точных повторов, число компаний рассчитано по уникальным идентификаторам.}
    \label{tab:data-snapshots}
    \begin{tabular}{@{}lrrr>{\raggedright\arraybackslash}p{5.7cm}@{}}
        \toprule
        Дата / версия & Строк & Компаний & Полей & Использование \\
        \midrule
        27.02.2023 & 4\,128 & 4\,128 & 16 & Проверка согласованности с июльским файлом \\
        13.07.2023 & 8\,410 & 4\,297 & 16 & Определение исходной группы и её характеристик \\
        01.08.2026 & 6\,110 & 6\,110 & 23 & Определение последующего статуса \\
        \bottomrule
    \end{tabular}
\end{table}

Февральскому срезу соответствует файл \path{2023-02-27-yc-companies.csv}, июльскому --- \path{2023-07-13-yc-companies.csv}. Поздняя версия представлена файлом \path{yc_companies.csv}.

Автор исследования скачал поздний файл со страницы Kaggle в середине августа 2026 года. Согласно его сведениям о странице источника, датасет обновляется ежемесячно первого числа. Поэтому 01.08.2026 принято как обозначение августовской версии, а не как подтверждённый момент сбора каждой строки. Номер версии Kaggle и точный день скачивания не сохранены. Утверждения о позднем состоянии относятся к содержимому этого файла; точный горизонт до календарного дня не устанавливается. Основание датировки записано в \path{data/source_provenance.json}.

Сведения получены из сторонних выгрузок, а не собраны непосредственно с сайта Y Combinator. Контрольные суммы по алгоритму SHA-256 рассчитываются в функции \texttt{main} файла \href{https://github.com/Kiraashado/course_work_startup_success_predictors/blob/main/src/data_preprocessing.py}{\texttt{src/data\_preprocessing.py}} и записываются в локальный отчёт \path{artifacts/preprocessed/cohort_report.json}. Контрольная сумма подтверждает идентичность файла, но не дату сбора. Полученные значения приведены в Таблице~\ref{tab:data-checksums}.

\begin{table}[!htbp]
    \centering
    \scriptsize
    \caption{Контрольные суммы исходных файлов по алгоритму SHA-256}
    \label{tab:data-checksums}
    \begin{tabular}{@{}lp{11.2cm}@{}}
        \toprule
        Файл & Контрольная сумма \\
        \midrule
        \path{2023-02-27-yc-companies.csv} & \texttt{ba602c1e186a6aa2407ed75f3081536e48519c8b6ed8f065e7a8273dc1cbe37c} \\
        \path{2023-07-13-yc-companies.csv} & \texttt{766d367b69cba93e370a868b3d62d5bdacf1dad1ba865085fd78cfa72241d942} \\
        \path{yc_companies.csv} & \texttt{c9336c008895c6f2bd6e53a47ce71ac509fe8fbb0bd72adde276dce05ed89ccc} \\
        \bottomrule
    \end{tabular}
\end{table}

\clearpage
\subsection{Поля исходных файлов}

Структура файлов 2023 года приведена в Таблице~\ref{tab:data-fields-2023}, структура позднего файла --- в Таблице~\ref{tab:data-fields-2026}. Если поле принимает произвольный текст или большое число различных категорий, в Таблицах указаны формат и несколько примеров; конечные наборы статусов и логических значений перечислены полностью. Пустые ячейки означают отсутствие сведений в источнике. Там, где поздняя выгрузка не содержит словаря полей, их описание ограничено названием и форматом записанных значений.

\begingroup
\footnotesize
\setstretch{1.1}
\begin{longtable}{@{}>{\RaggedRight\arraybackslash}p{3.6cm}>{\RaggedRight\arraybackslash}p{4.8cm}>{\RaggedRight\arraybackslash}p{7.7cm}@{}}
\caption{Поля февральской и июльской выгрузок 2023 года. Названия и примеры значений сохранены в форме исходных файлов.}
\label{tab:data-fields-2023}\\
\toprule
Поле & Возможные значения & Содержание \\
\midrule
\endfirsthead
\multicolumn{3}{@{}l}{Продолжение Таблицы~\thetable}\\
\toprule
Поле & Возможные значения & Содержание \\
\midrule
\endhead
\bottomrule
\endlastfoot
\path{company_id} & Положительное целое число & Идентификатор компании в каталоге; служит ключом при сопоставлении выгрузок. \\
\path{company_name} & Текст, например \texttt{moonrepo} & Название компании; совпадение названий не заменяет совпадения идентификаторов. \\
\path{short_description} & Короткий текст либо пропуск & Однострочное описание деятельности компании. \\
\path{long_description} & Текст либо пропуск & Развёрнутое описание деятельности, если оно заполнено в карточке. \\
\path{batch} & Код набора, например \texttt{W23}, \texttt{S22} или \texttt{IK12} & Обозначение набора Y Combinator. Буквы \texttt{W} и \texttt{S} обозначают зимний и летний наборы; \texttt{IK12} не имеет стандартного сезонного обозначения. \\
\path{status} & \texttt{Active}, \texttt{Inactive}, \texttt{Acquired}, \texttt{Public} & Статус компании на дату соответствующей выгрузки. \\
\path{tags} & Список строк, например \texttt{['SaaS', 'B2B']}, либо пустой список & Тематические метки карточки. Одна компания может иметь несколько меток. \\
\path{location} & Текст, например \texttt{Portland, OR}, либо пропуск & Местоположение в исходной записи без выделения страны и города по единому правилу. \\
\path{country} & Краткий код страны, например \texttt{US} или \texttt{IN}, либо пропуск & Страна, выделенная в источнике в отдельное поле. \\
\path{year_founded} & Год основания либо пропуск & Указанный в карточке год основания; точная дата не содержится в этом поле. \\
\path{num_founders} & Целое число от 0 до 6 & Число основателей в источнике. Ноль при подготовке считается неизвестным значением. \\
\path{founders_names} & Список имён, в том числе пустой & Имена основателей в строковом представлении списка. \\
\path{team_size} & Неотрицательное целое число либо пропуск & Численность команды по карточке компании; ноль сохраняется как отдельное наблюдаемое значение. \\
\path{website} & Ссылка на сайт либо пропуск & Адрес сайта компании. \\
\path{cb_url} & Ссылка на Crunchbase либо пропуск & Адрес карточки компании в каталоге Crunchbase. \\
\path{linkedin_url} & Ссылка на LinkedIn либо пропуск & Адрес страницы компании в LinkedIn. \\
\end{longtable}
\endgroup

В поздней выгрузке отсутствуют отдельные поля страны, года основания, числа и имён основателей. Эти сведения не восстанавливались по поздним адресам или другим полям: признаки для анализа исходного состояния взяты из июльского файла.

Одноимённые по смыслу поля приведены к общей схеме: \path{company_id} раннего файла и \path{id} позднего получили название \path{id}. Поля \path{launched_at}, \path{industries}, \path{regions}, \path{app_video_public}, \path{demo_day_video_public}, \path{app_answers} и \path{question_answers} не вошли в очищенный поздний срез. Из него к исходным записям присоединён только статус; остальные поздние характеристики не использованы как признаки прогноза.

\begingroup
\footnotesize
\setstretch{1.1}
\begin{longtable}{@{}>{\RaggedRight\arraybackslash}p{3.6cm}>{\RaggedRight\arraybackslash}p{4.8cm}>{\RaggedRight\arraybackslash}p{7.7cm}@{}}
\caption{Все поля выгрузки августовской версии 2026 года. Для перечислений с множеством значений приведены формат и примеры.}
\label{tab:data-fields-2026}\\
\toprule
Поле & Возможные значения & Содержание \\
\midrule
\endfirsthead
\multicolumn{3}{@{}l}{Продолжение Таблицы~\thetable}\\
\toprule
Поле & Возможные значения & Содержание \\
\midrule
\endhead
\bottomrule
\endlastfoot
\path{id} & Положительное целое число & Идентификатор компании; сопоставляется с полем \path{company_id} ранней выгрузки. \\
\path{name} & Текст & Название компании в позднем каталоге. \\
\path{website} & Ссылка на сайт либо пропуск & Адрес сайта компании на дату поздней выгрузки. \\
\path{all_locations} & Перечень мест текстом либо пропуск & Все местоположения, указанные в карточке; несколько мест могут быть разделены точкой с запятой. \\
\path{long_description} & Текст либо пропуск & Развёрнутое описание компании в поздней карточке. \\
\path{one_liner} & Короткий текст либо пропуск & Однострочное описание компании; по смыслу соответствует раннему \path{short_description}. \\
\path{team_size} & Неотрицательное целое число либо пропуск & Численность команды в поздней карточке. Для исходных признаков это значение не используется. \\
\path{industry} & Девять категорий: \texttt{B2B}, \texttt{Consumer}, \texttt{Education}, \texttt{Fintech}, \texttt{Government}, \texttt{Healthcare}, \texttt{Industrials}, \texttt{Real Estate and Construction}, \texttt{Unspecified} & Широкая категория деятельности в классификации позднего источника. \\
\path{subindustry} & Категория, например \texttt{B2B -> Infrastructure} & Более подробное направление деятельности; в файле встречаются 59 обозначений. \\
\path{launched_at} & Дата и время, например \texttt{2011-11-23 10:52:03} & Метка времени запуска, указанная источником. Она не приравнивается к году основания компании. \\
\path{tags} & Список строк либо пустой список & Тематические метки в поздней карточке. Они могут отличаться от меток 2023 года. \\
\path{top_company} & \texttt{True}, \texttt{False} либо пропуск & Отметка компании как одной из выделенных в каталоге. \\
\path{isHiring} & \texttt{True} или \texttt{False} & Признак открытого найма в поздней карточке. \\
\path{nonprofit} & \texttt{True} или \texttt{False} & Признак некоммерческой организации. \\
\path{batch} & Сезон и год, например \texttt{Winter 2022}; также \texttt{Unspecified} & Обозначение набора Y Combinator в формате позднего источника. В файле есть весенние и осенние наборы, а также одна запись с годом 2027. \\
\path{status} & \texttt{Active}, \texttt{Inactive}, \texttt{Acquired}, \texttt{Public} & Статус компании на дату поздней выгрузки; единственное содержательное поле, переносимое в исследуемую группу. \\
\path{industries} & Список строк, например \texttt{['B2B', 'Infrastructure']} & Перечень уровней отраслевой классификации в исходной карточке. \\
\path{regions} & Список строк, например \texttt{['Europe', 'Remote']} & Географические и связанные с удалённой работой метки источника; это не отдельное поле страны. \\
\path{stage} & \texttt{Early} или \texttt{Growth} & Категория стадии развития в позднем каталоге. \\
\path{app_video_public} & \texttt{True} или \texttt{False} & Признак публичной доступности видеозаписи из заявки в Y Combinator. \\
\path{demo_day_video_public} & \texttt{True} или \texttt{False} & Признак публичной доступности видеозаписи выступления на презентационном дне. \\
\path{app_answers} & \texttt{True}, \texttt{False} либо пропуск & Признак, связанный с ответами из заявки. В 6\,033 из 6\,110 строк значение отсутствует; смысл пропуска источник не разъясняет. \\
\path{question_answers} & \texttt{True} или \texttt{False} & Признак публичной доступности ответов компании на вопросы в каталоге. \\
\end{longtable}
\endgroup

\subsection{Формирование исследуемой группы и производные поля}

Единица наблюдения --- компания с уникальным идентификатором. Исходная группа состоит из 3\,110 компаний со статусом \texttt{Active} в июльской выгрузке. Февральский файл использован для проверки согласованности записей, но не для заполнения июльских пропусков. Производные поля рассчитаны после удаления повторов и очистки исходных значений.

В Таблице~\ref{tab:data-derived} перечислены производные характеристики, технические поля и целевая переменная. Технические поля предназначены для учёта происхождения записей и сопоставления срезов; в прогноз они не включены.

\begingroup
\footnotesize
\setstretch{1.1}
\begin{longtable}{@{}>{\RaggedRight\arraybackslash}p{3.6cm}>{\RaggedRight\arraybackslash}p{4.8cm}>{\RaggedRight\arraybackslash}p{7.7cm}@{}}
\caption{Производные и технические поля подготовленных данных. Расчёт характеристик компании выполняется отдельно внутри каждого среза; в прогноз поступают только характеристики июля 2023 года.}
\label{tab:data-derived}\\
\toprule
Поле & Возможные значения & Способ получения \\
\midrule
\endfirsthead
\multicolumn{3}{@{}l}{Продолжение Таблицы~\thetable}\\
\toprule
Поле & Возможные значения & Способ получения \\
\midrule
\endhead
\bottomrule
\endlastfoot
\path{tags} & Список различных меток; возможен пустой список & Строковое представление разбирается в список; пробелы по краям меток удаляются, повторы исключаются. \\
\path{tag_count} & Неотрицательное целое число & Число различных меток после разбора \path{tags}; пустому списку соответствует ноль. \\
\path{description_length} & Неотрицательное целое число либо пропуск & Число символов в очищенном кратком описании. При отсутствии описания длина остаётся неизвестной. \\
\path{company_age} & Неотрицательное число лет либо пропуск & Год среза минус \path{year_founded}. Это разность календарных лет, а не возраст в полных годах. \\
\path{batch_year} & Год набора либо пропуск & Год выделяется из \path{batch}. Если он позже года среза, числовое значение признаётся недопустимым. \\
\path{batch_season} & \texttt{Winter}, \texttt{Summer}, \texttt{Spring}, \texttt{Fall} либо пропуск & Сезон выделяется из стандартного кода набора. Для \texttt{IK12} сезон не назначается. \\
\path{regular_batch} & \texttt{True} или \texttt{False} & Показывает, распознан ли обычный код с сезоном и годом. Для \texttt{IK12} и нераспознанных обозначений значение ложно. \\
\path{future_batch} & \texttt{True} или \texttt{False} & Показывает, что извлечённый год набора позже года соответствующего среза. Исходный код при этом сохраняется, а \path{batch_year} становится пропуском. \\
\path{snapshot_date} & Дата соответствующей выгрузки & Постоянная дата, присвоенная всем строкам очищенного среза. \\
\path{source} & \path{yc_2023_feb}, \path{yc_2023_jul} или \path{yc_2026} & Обозначение файла-источника, присвоенное при загрузке. \\
\path{status_baseline} & \texttt{active} & Июльский статус после отбора исходной группы; исходное поле \path{status} переименовано. \\
\path{status_followup} & \texttt{active}, \texttt{inactive}, \texttt{acquired}, \texttt{public}, \texttt{unknown} & Статус сопоставленной компании из позднего файла. При отсутствии записи проставляется \texttt{unknown}. \\
\path{followup_found} & \texttt{True} или \texttt{False} & Результат сопоставления идентификатора с поздним файлом. \\
\path{outcome_known} & \texttt{True} или \texttt{False} & Показывает, входит ли поздний статус в четыре распознаваемые категории. \\
\path{success_binary} & 1, 0 либо пропуск & Единица присваивается статусам \texttt{acquired} и \texttt{public}; ноль --- статусам \texttt{active} и \texttt{inactive}. При неизвестном статусе значение отсутствует. \\
\path{followup_date} & \texttt{2026-08-01} & Обозначение августовской версии; не дата обновления каждой карточки. \\
\end{longtable}
\endgroup

При очистке пустые строки заменены пропусками, статусы приведены к строчным буквам. Коды стран записаны заглавными буквами: \texttt{USA} и \texttt{UNITED STATES} заменены на \texttt{US}, \texttt{UK} и \texttt{UNITED KINGDOM} --- на \texttt{GB}. Нулевое число основателей переведено в пропуск, нулевой размер команды сохранён. Отрицательные и дробные значения размера команды, числа основателей и года основания признаны недопустимыми. Год основания также переведён в пропуск, если он равен нулю или превышает год выгрузки.

В прогнозном эксперименте размер команды представлен как $\ln(1+x)$, где $x$ обозначает число человек. Неизвестный сезон выделен в отдельную категорию. При обучении логистической регрессии пропуски числовых признаков отмечены отдельными индикаторами, а год набора и возраст представлены сплайнами. Заполнение числовых пропусков и оценка параметров преобразования проводились только на обучающих частях. Подробности приведены в разделе~\ref{sec:modeling-methods}.

\subsection{Последующий статус и границы наблюдения}

Исходные и поздние записи сопоставлены по уникальному идентификатору с сохранением всех компаний июльской группы. У 3\,046 компаний поздний статус установлен; 64 компании в файле августовской версии 2026 года не найдены и сохранены с неизвестным исходом. Компании, присутствующие только в поздней выгрузке, в исследуемую группу не вошли.

Положительная метка присвоена поглощённым и публичным компаниям, отрицательная --- активным и неактивным в августовской версии 2026 года. При неизвестном позднем статусе метка отсутствует. Отрицательный класс означает лишь отсутствие выбранного корпоративного события на конечную дату; финансовый результат компании и дата события по этим данным не определяются.

\subsection{Проверки и воспроизводимость}

Удалены только полностью совпадающие строки. Отсутствующий либо некорректный идентификатор и несовпадающие записи с одним идентификатором считались ошибками, препятствующими сопоставлению; в использованных файлах таких конфликтов после удаления повторов не обнаружено. Коду \texttt{IK12} сезон не присвоен. При разборе полных обозначений наборов \texttt{Spring} отделён от \texttt{Summer}.

Сохранены очищенные срезы, исследуемая группа и машиночитаемый отчёт о подготовке. Отчёт фиксирует размеры и контрольные суммы файлов, пропуски, удалённые повторы, недопустимые числовые значения и результаты сопоставления. Для разведочного анализа использована сохранённая группа компаний; сводные Таблицы, лежащие в основе Рисунков, также записаны в отдельные файлы.
\clearpage
